% Final replacement for Sections 3--5 of the DR-DAD manuscript.
% This is a manuscript fragment: input it after the paper preamble.
% Required packages: amsmath, amssymb, booktabs, graphicx, algorithm,
% algpseudocode, and threeparttable (only if the host manuscript uses it).
%
% Reporting scope:
% - Noto 2024 is the primary empirical case.
% - Turkey 2023 and Nepal 2015 are supplementary cross-case checks.
% - Five-link policies are optimal only over the declared finite grid.
% - The outcome is modeled loss / people without timely access, not deaths.

\section{Model formulation}
\label{sec:model}

\subsection{Sets, decisions, and modeled outcome}
\label{sec:notation}

Let \(\mathcal R\) be the set of demand zones, \(\mathcal K\) the set of
emergency-response centers, and \(\mathcal L\) the set of candidate road
links. A zone index \(r\) may represent the composite index \(rl\) used in
the original geographic notation. A road-failure state \(s\in\mathcal S\)
is identified by the subset \(F(s)\subseteq\mathcal L\) of failed candidate
links. The primary empirical analysis uses the complete no-tail state space,
so \(\lvert\mathcal S\rvert=2^{\lvert\mathcal L\rvert}\). A reduced state
space with a residual tail state is used only in computational diagnostics;
it is not used to identify the primary empirical policy.

The first-stage decisions are building-renovation fractions
\(z_r\in[0,1]\), road-retrofit fractions \(y_\ell\in[0,1]\), and additional
response capacity \(w_k\geq0\). Their feasible set is
\begin{align}
\mathcal X
  =\bigg\{(z,w,y):\;&
     \sum_{r\in\mathcal R} C_r^Z z_r \leq B_Z,\quad
     \sum_{k\in\mathcal K} C_k^X w_k \leq B_X,\quad
     \sum_{\ell\in\mathcal L} C_\ell^Y y_\ell \leq B_Y,
     \nonumber\\[-2mm]
   &0\leq z_r\leq1,\quad w_k\geq0,\quad 0\leq y_\ell\leq1
   \bigg\}.
\label{eq:feasible-set}
\end{align}
Here \(C_r^Z\), \(C_\ell^Y\), and \(C_k^X\) are renovation, road-retrofit,
and capacity-expansion cost coefficients, respectively; \(B_Z\), \(B_Y\),
and \(B_X\) are sector-specific budgets.

For zone \(r\), let \(P_r\) be population and \(q_r\) the severe-damage or
collapse proxy. The baseline at-risk population is
\begin{equation}
  A_r=P_rq_r,
  \qquad
  D(z)=\sum_{r\in\mathcal R} A_r(1-z_r).
\label{eq:at-risk-population}
\end{equation}
The framework intentionally models a specified acute-response component of
this population rather than treating \(D(z)\) as a forecast of fatalities.
Let \(\eta_r\in[0,1]\) be the fraction requiring time-sensitive response and
\(\delta_r\in[0,1]\) the fraction assigned to immediate modeled loss, with
\(\eta_r+\delta_r\leq1\). After renovation, time-sensitive demand and
immediate modeled loss are
\begin{equation}
  d_r(z)=\eta_r A_r(1-z_r),
  \qquad
  \iota_r(z)=\delta_r A_r(1-z_r).
\label{eq:modeled-demand}
\end{equation}
The remaining part of \(A_r(1-z_r)\) is outside the modeled acute-response
outcome and is not interpreted as either a survivor or a death.

After a road-failure state is observed, the responder dispatches
\(x_{kr}^s\geq0\) units from center \(k\) to zone \(r\). Let
\(g_{kr}^s(y)\in[0,1]\) be the response benefit per dispatched unit. The
primary Noto specification uses a timely-access indicator,
\begin{equation}
 g_{kr}^s(y)=\mathbf{1}\{\tau_{kr}^s(y)\leq T\},
\label{eq:timely-access}
\end{equation}
where \(T\) is a declared response-time threshold. A calibrated continuous
response function can be substituted for \eqref{eq:timely-access}; doing so
changes the outcome definition and requires an independently justified
calibration. The statewise recourse problem is
\begin{align}
 Q_s(z,w;y)=\max_{x^s\geq0}\quad&
     \sum_{k\in\mathcal K}\sum_{r\in\mathcal R}
       g_{kr}^s(y)x_{kr}^s \label{eq:recourse-primal}\\
 \text{s.t.}\quad&
     \sum_{r\in\mathcal R}x_{kr}^s\leq w_k^0+w_k,
       &&k\in\mathcal K, \nonumber\\
 &   \sum_{k\in\mathcal K}x_{kr}^s\leq d_r(z),
       &&r\in\mathcal R. \nonumber
\end{align}
Thus \(Q_s\) is the maximum number of time-sensitive people receiving
timely access under the scenario interpretation. The modeled state loss is
\begin{equation}
 \gamma_s(z,w;y)
   =\sum_{r\in\mathcal R}\left[d_r(z)+\iota_r(z)\right]-Q_s(z,w;y).
\label{eq:state-loss}
\end{equation}
It is nonnegative under \eqref{eq:recourse-primal}. Lower objective values
therefore indicate fewer modeled losses, not a calibrated prediction of
observed mortality.

\subsection{Road reliability, conditional performance, and nominal states}
\label{sec:road-model}

For link \(\ell\), let \(\Phi_\ell\) be the baseline failure probability and
\(\underline\Phi_\ell\) a residual failure floor satisfying
\(0\leq\underline\Phi_\ell\leq\Phi_\ell\leq1\). Road retrofit follows
\begin{equation}
 \phi_\ell(y_\ell)
   =\underline\Phi_\ell+
       \left(\Phi_\ell-\underline\Phi_\ell\right)(1-y_\ell)
   =\Phi_\ell-\left(\Phi_\ell-\underline\Phi_\ell\right)y_\ell.
\label{eq:residual-reliability}
\end{equation}
The residual floor distinguishes strengthening from complete replacement:
full retrofit reduces, but need not eliminate, modeled failure risk. The
zero-floor case is retained as a limiting engineering assumption only when
it is substantively justified.

Road retrofit also affects conditional access after a failure. Let
\(\overline\tau_{kr}\) be the intact-route travel time and let
\(\Delta\tau_{kr,\ell}\geq0\) be the observed or constructed disruption
penalty attributable to failed link \(\ell\) on the route from \(k\) to
\(r\). Let \(\alpha_\ell\in[0,1]\) denote the fraction by which retrofit
reduces that penalty. We use
\begin{equation}
 \tau_{kr}^s(y)
 =\overline\tau_{kr}+
   \sum_{\ell\in F(s)}
      \Delta\tau_{kr,\ell}(1-\alpha_\ell y_\ell).
\label{eq:conditional-delay}
\end{equation}
Equation \eqref{eq:conditional-delay} makes the physical mechanism explicit:
road retrofit affects both the likelihood of each failure state and the
severity of disrupted access conditional on that state. It is an aggregated
corridor representation, not a dynamic traffic-assignment or rerouting
model.

Conditional on \(y\), link failures are modeled as independent for scenario
generation. Hence the nominal probability of state \(s\) is
\begin{equation}
 \pi_s(y)=
   \prod_{\ell\in F(s)}\phi_\ell(y_\ell)
   \prod_{\ell\in\mathcal L\setminus F(s)}
       \left[1-\phi_\ell(y_\ell)\right].
\label{eq:nominal-probability}
\end{equation}
Independence is a transparent scenario assumption, not a claim that
earthquake-induced road failures are empirically independent. Dependence is
stress-tested through the ambiguity set and, in a supplementary sensitivity,
through moment restrictions.

\subsection{Capped decision-dependent ambiguity}
\label{sec:capped-tv}

For a fixed retrofit vector \(y\), nature selects a distribution \(p\) in
the capped total-variation set
\begin{equation}
\begin{split}
 \mathcal P_{\kappa,\rho}(y)=\bigg\{p\in\mathbb R_+^{\lvert\mathcal S\rvert}:
 &\ \sum_{s\in\mathcal S}p_s=1,\quad
 \frac{1}{2}\sum_{s\in\mathcal S}\lvert p_s-\pi_s(y)\rvert\leq\rho,
 \\[-1mm]
 &\ p_s\leq\kappa\pi_s(y),\quad s\in\mathcal S\bigg\},
 \qquad \kappa\geq1.
\end{split}
\label{eq:capped-tv}
\end{equation}
The density cap repairs the physical inconsistency of unrestricted
total-variation ambiguity: when \(\pi_s(y)=0\), it forces \(p_s=0\), so
nature cannot assign mass to a state that the nominal model declares
impossible. For any event \(T\subseteq\mathcal S\),
\begin{equation}
 p(T)-\pi(T)
 \leq
 \min\left\{\rho,\;(\kappa-1)\pi(T),\;1-\pi(T)\right\}.
\label{eq:rare-state-bound}
\end{equation}
Thus a rare nominal event cannot receive arbitrary probability merely because
the total-variation radius is positive.

When residual floors in \eqref{eq:residual-reliability} are strictly positive,
all modeled no-tail states generally have positive nominal mass. In that
case \eqref{eq:capped-tv} is more precisely a bounded likelihood-ratio
restriction than a literal state-elimination rule. Literal support
preservation is obtained in the zero-floor limit. Both interpretations are
useful: residual floors encode imperfect engineering performance, while the
cap prevents implausibly large probability amplification.

For nonnegative state losses, the robust expectation obeys
\begin{equation}
\begin{split}
\sum_s\pi_s(y)\gamma_s
&\leq
 \max_{p\in\mathcal P_{\kappa,\rho}(y)}
     \sum_s p_s\gamma_s\\
&\leq
 \min\left\{
   \sum_s\pi_s(y)\gamma_s+
       \rho\left(\max_s\gamma_s-\min_s\gamma_s\right),\
   \kappa\sum_s\pi_s(y)\gamma_s
 \right\}.
\end{split}
\label{eq:ambiguity-premium-bound}
\end{equation}
Equation \eqref{eq:ambiguity-premium-bound} is used as an audit check, not
as a replacement for solving the adversarial linear program.

\paragraph{Optional moment-aware sensitivity.}
The primary specification is \eqref{eq:capped-tv}. To test whether the
total-variation ball permits implausible changes in marginal failure rates
or failure dependence, we also consider an optional intersection
\(\mathcal P_{\kappa,\rho}^{\mathrm{mom}}(y)\). Let
\(f_{\ell s}=\mathbf{1}\{\ell\in F(s)\}\),
\(\mu_\ell(y)=\sum_s\pi_s(y)f_{\ell s}\), and
\(N_s=\sum_\ell f_{\ell s}\). The sensitivity adds linear, fixed-\(y\)
bands such as
\begin{align}
\left\lvert\sum_s p_s f_{\ell s}-\mu_\ell(y)\right\rvert
 &\leq a_1+r_1\mu_\ell(y), &&\ell\in\mathcal L, \label{eq:moment-marginal}\\
\left\lvert\sum_s p_sN_s-\sum_s\pi_s(y)N_s\right\rvert
 &\leq\varepsilon_\mu, \label{eq:moment-mean}\\
\left\lvert\sum_s p_s
  \left[N_s-\sum_t\pi_t(y)N_t\right]^2-M_2(y)\right\rvert
 &\leq\varepsilon_2(y), \label{eq:moment-second}
\end{align}
where \(M_2(y)\) is the nominal fixed-center second moment. Analogous bands
can be imposed on \(f_{\ell s}f_{\ell' s}\). These constraints are linear in
\(p\) for fixed \(y\), so they retain the exact fixed-policy oracle. Their
tolerances are scenario parameters unless supported by a hazard ensemble;
they are not estimated moments in the present cases.

\subsection{DR-DAD problem and interpretation}
\label{sec:dad-problem}

The decision-dependent distributionally robust disaster-risk-reduction
problem is
\begin{equation}
 \min_{(z,w,y)\in\mathcal X}
 \quad V_\rho(z,w,y)
 :=
 \max_{p\in\mathcal P_{\kappa,\rho}(y)}
 \sum_{s\in\mathcal S}p_s\gamma_s(z,w;y).
\label{eq:dr-dad}
\end{equation}
The decision-dependent feature has two channels. First, \(y\) changes
\(\pi(y)\) and therefore the ambiguity set in \eqref{eq:capped-tv}. Second,
\(y\) changes conditional travel times through \eqref{eq:conditional-delay}.
The latter channel was absent from the earlier static implementation and is
essential for interpreting road hardening as an access intervention rather
than only a probability adjustment.

The framework is designed to identify a defensible policy under declared
hazard, engineering, capacity, and ambiguity assumptions. It does not claim
that observed post-earthquake damage alone identifies \(\Phi_\ell\),
\(\underline\Phi_\ell\), \(\alpha_\ell\), costs, budgets, capacity
operability, \(\rho\), or \(\kappa\); those inputs are reported separately
as constructed or scenario-calibrated quantities.


\section{Solution methodology}
\label{sec:solution}

\subsection{Exact fixed-road evaluation}
\label{sec:fixed-y}

For a fixed road plan \(y\), the nominal probabilities, conditional travel
times, and ambiguity set are fixed. Define
\begin{equation}
 \Theta_\rho(y)
 =
 \min_{(z,w):(z,w,y)\in\mathcal X}
 \max_{p\in\mathcal P_{\kappa,\rho}(y)}
 \sum_s p_s\gamma_s(z,w;y).
\label{eq:fixed-y-value}
\end{equation}
This problem is solved by finite cut generation. The dual of the
state-\(s\) recourse problem \eqref{eq:recourse-primal} is
\begin{align}
 Q_s(z,w;y)
 =\min_{a^s,b^s\geq0}\quad&
   \sum_{k\in\mathcal K}(w_k^0+w_k)a_k^s+
   \sum_{r\in\mathcal R}d_r(z)b_r^s \label{eq:recourse-dual}\\
 \text{s.t.}\quad&
   a_k^s+b_r^s\geq g_{kr}^s(y),
      &&k\in\mathcal K,\ r\in\mathcal R. \nonumber
\end{align}
Each dual extreme point \(c=(a^s,b^s)\) yields the valid affine loss cut
\begin{equation}
 \theta_s\geq
 H_{s,c}(z,w;y)
 :=
 \sum_r[d_r(z)+\iota_r(z)]
 -\sum_k(w_k^0+w_k)a_k^s-\sum_r d_r(z)b_r^s .
\label{eq:recourse-cut}
\end{equation}
The master problem contains all currently generated cuts
\(\theta_s\geq H_{s,c}\), the feasible set \eqref{eq:feasible-set}, and the
finite linear adversarial problem over \(p\). If a newly evaluated statewise
dispatch violates its current loss approximation, its dual extreme point is
added. Otherwise, the capped-TV adversary is solved exactly as a linear
program. The optional moment restrictions
\eqref{eq:moment-marginal}--\eqref{eq:moment-second} add only linear rows to
that fixed-\(y\) adversarial problem.

\begin{algorithm}[t]
\caption{Fixed-road oracle for \(\Theta_\rho(y)\)}
\label{alg:fixed-y}
\begin{algorithmic}[1]
\Require Feasible \(y\), tolerance \(\epsilon_{\mathrm{in}}\)
\State Construct \(\pi(y)\), \(\tau^s(y)\), \(g^s(y)\), and
       \(\mathcal P_{\kappa,\rho}(y)\)
\State Initialize each state cut set \(\mathcal C_s\) with a feasible
       recourse dual point
\Repeat
  \State Solve the master problem in \((z,w,\theta,p)\) using
         \(\{\theta_s\geq H_{s,c}:c\in\mathcal C_s\}\)
  \ForAll{\(s\in\mathcal S\)}
    \State Solve \eqref{eq:recourse-primal} at the candidate \((z,w,y)\)
    \If{the state loss violates \(\theta_s\) by more than
        \(\epsilon_{\mathrm{in}}\)}
      \State Add the associated dual cut \eqref{eq:recourse-cut} to
             \(\mathcal C_s\)
    \EndIf
  \EndFor
\Until{no recourse cut is violated}
\State Return the fixed-\(y\) objective, \((z,w)\), state losses, and
       worst-case distribution
\end{algorithmic}
\end{algorithm}

Because there are finitely many states and the recourse dual is polyhedral,
Algorithm~\ref{alg:fixed-y} is an exact fixed-policy evaluation method up
to the stated numerical tolerance. All reported policies are accompanied by
probability-sum, nonnegativity, budget, and inner-gap checks.

\subsection{Exhaustive search over an implementable road grid}
\label{sec:grid-search}

The final empirical solver is exhaustive enumeration over a predeclared,
implementable grid rather than an unsupported continuous global claim. For
the five-corridor case we use
\begin{equation}
 \mathcal G
 =
 \left\{0,0.25,0.50,0.75,1\right\}^{\lvert\mathcal L\rvert}
 \cap
 \left\{y:\sum_{\ell\in\mathcal L}C_\ell^Yy_\ell\leq B_Y\right\}.
\label{eq:road-grid}
\end{equation}
Every \(y\in\mathcal G\) is evaluated by Algorithm~\ref{alg:fixed-y}; the
selected solution
\begin{equation}
  \widehat y_\rho\in\arg\min_{y\in\mathcal G}\Theta_\rho(y)
\label{eq:grid-policy}
\end{equation}
is therefore globally best over \(\mathcal G\), but not certified as
globally optimal on the continuous box \([0,1]^{\lvert\mathcal L\rvert}\).
This distinction is retained in all tables and claims.

We report policy adaptation using complete plans, not only road vectors.
Let \((\widehat z_0,\widehat w_0,\widehat y_0)\) be the policy selected at
\(\rho=0\), and let
\((\widehat z_\rho,\widehat w_\rho,\widehat y_\rho)\) be the policy selected
at radius \(\rho\). The value of adapting the complete plan is
\begin{equation}
 \Delta_\rho^{\mathrm{all}}
 =
 V_\rho(\widehat z_0,\widehat w_0,\widehat y_0)
 -
 V_\rho(\widehat z_\rho,\widehat w_\rho,\widehat y_\rho).
\label{eq:complete-adaptation}
\end{equation}
This is evaluated with the whole \(\rho=0\) plan fixed in the first term; it
does not re-optimize \(z\) or \(w\). A zero
\(\Delta_\rho^{\mathrm{all}}\) means that ambiguity changes evaluation but
provides no value from changing the selected complete grid policy. For
policy identification, we also report the number of near-optimal road
portfolios
\begin{equation}
 N_\varepsilon(\rho)
 =
 \left\lvert
   \left\{y\in\mathcal G:
   \frac{\Theta_\rho(y)-\min_{\tilde y\in\mathcal G}
   \Theta_\rho(\tilde y)}
   {\min_{\tilde y\in\mathcal G}\Theta_\rho(\tilde y)}
   \leq\varepsilon
   \right\}
 \right\rvert .
\label{eq:near-optimal-count}
\end{equation}
Large values of \(N_\varepsilon(\rho)\) identify a flat discrete policy
landscape rather than a meaningful ambiguity-induced policy switch.

\subsection{Continuous spatial branch-and-bound: restricted diagnostic}
\label{sec:sbb}

The former continuous spatial branch-and-bound (SBB) formulation is retained
only as a static probability-only diagnostic. In that restricted model,
\(y\) enters \(\pi_s(y)\) but not \(\gamma_s\), so product-tree variables and
McCormick envelopes can relax the multilinear state probabilities. A
nonnegative lower bound is imposed explicitly:
\begin{equation}
 LB(\mathcal B)\leftarrow\max\{0,LB_{\mathrm{LP}}(\mathcal B)\}.
\label{eq:valid-lb}
\end{equation}
The final model \eqref{eq:conditional-delay} has decision-dependent
conditional travel times; the current SBB relaxation deliberately rejects
that model rather than returning a bound that omits its physical channel.
Likewise, it is not used for the optional moment-aware ambiguity extension.

The five-link static root relaxation is too weak for a continuous empirical
certificate: its valid lower bound is zero. A corner-bound ablation reduced
the number of LP variables but did not improve this lower bound. We
therefore do not report a five-link continuous SBB optimum, do not use it to
claim that road policies are globally optimal in the continuous model, and
do not present corner bounds as an exact replacement for the multilinear
probability graph. Small nested static instances are useful code
verification tests only when their incumbent-bound gap closes.

\begin{table}[t]
\centering
\caption{Static SBB root-relaxation diagnostic.}
\label{tab:sbb-root-diagnostic}
\small
\begin{tabular}{llrrr}
\toprule
Candidate links & Relaxation diagnostic & Raw root LB & Valid root LB & LP variables\\
\midrule
2 & Product tree & \(-60{,}684.736\) & 0 & 38\\
3 & Product tree & \(-141{,}342.025\) & 0 & 63\\
4 & Product tree & \(-172{,}760.166\) & 0 & 120\\
5 & Product tree & \(-184{,}026.202\) & 0 & 249\\
5 & Corner-bound ablation & \(-184{,}026.202\) & 0 & 121\\
\bottomrule
\end{tabular}
\par\smallskip
\footnotesize\textit{Notes.} These are static probability-only diagnostics.
They are not lower bounds for the final conditional-performance model in
\eqref{eq:conditional-delay}. The valid nonnegative lower bound is
uninformative for the five-link continuous problem.
\end{table}

\subsection{Restart-safe execution and audit trail}
\label{sec:computational-audit}

The enumeration writes an atomic checkpoint after each completed road
portfolio. Each record contains the grid index, \(y\), the fixed-\(y\)
solution, objective, numerical checks, and a timestamp. On restart, the
driver validates the checkpoint and skips completed portfolios; a detached
worker and append-only progress log allow monitoring without keeping the
interactive application open. Thus an application interruption does not
invalidate completed fixed-policy evaluations. Source manifests, parameter
files, run identifiers, and result tables are retained with the computation
to make the reported grid certificate auditable.


\section{Empirical implementation and results}
\label{sec:results}

\subsection{Case-study role and data-to-model mapping}
\label{sec:empirical-data}

The 2024 Noto Peninsula earthquake is the primary empirical implementation
because public sources provide an unusually direct record of the access
mechanism: municipality population and households, post-event dwelling
damage, official intercity corridor restoration information and travel-time
observations, and pre-event hospital beds. This combination supports a
transparent calibration of exposure, damaged dwellings, baseline access, and
observed access disruption. It does not identify all behavioral or
engineering parameters in \eqref{eq:dr-dad}.

\begin{table}[t]
\centering
\caption{Noto source layers and status in the DR-DAD implementation.}
\label{tab:noto-data-status}
\small
\resizebox{\textwidth}{!}{%
\begin{tabular}{p{3.1cm}p{4.2cm}p{2.9cm}p{5.0cm}}
\toprule
Model layer & Public source & Status & Use and limitation\\
\midrule
Population and households &
Ishikawa Prefecture Statistics (2022) &
Observed &
\(P_r\); household denominator used in the damage proxy.\\
Dwelling damage &
Ishikawa damage report No.~162 (1 October 2024) &
Observed &
Fully destroyed dwellings; used to construct \(q_r\), not person-level
casualties.\\
Corridor geometry, disruption, and recovery &
MLIT Noto road-restoration GIS and travel-time records &
Observed &
Candidate corridors, intact and disrupted route times, and
\(\Delta\tau_{kr,\ell}\).\\
Hospital beds &
MHLW FY2023 Hospital Function Report &
Observed &
Anchors \(w_k^0\); operational share and throughput conversion remain
scenario-calibrated.\\
Failure probabilities, residual floors, and delay-reduction factors &
Road delay ratios, restoration evidence, and declared engineering settings &
Constructed / scenario-calibrated &
\(\Phi_\ell,\underline\Phi_\ell,\alpha_\ell\); not repeated-event failure
frequencies or measured retrofit effects.\\
Costs, budgets, \(\eta,\delta,\rho,\kappa\) &
Declared scenario design and sensitivity analysis &
Scenario-calibrated &
Cost coefficients, sector budgets, modeled response fraction, and
ambiguity controls; not observed data.\\
\bottomrule
\end{tabular}}
\end{table}

The five demand zones are Nanao, Wajima, Suzu, Anamizu, and Noto. Kanazawa
is retained as an external high-capacity center, yielding six modeled
centers. The five candidate corridors are Kanazawa--Nanao,
Nanao--Anamizu, Anamizu--Wajima, Anamizu--Noto, and Anamizu--Suzu. Their
complete no-tail state space has \(2^5=32\) states.

For each municipality, the severe-damage proxy is
\begin{equation}
 q_r=\frac{\text{fully destroyed dwellings in }r}
             {\text{pre-event households in }r}.
\label{eq:noto-damage-proxy}
\end{equation}
This assumes the same average household size in destroyed and non-destroyed
households. The five zones contain 106,022 residents and have
\(D(0)=12{,}108.462\) baseline at-risk persons under
\eqref{eq:noto-damage-proxy}. The estimated damage fractions are 0.0250,
0.2375, 0.3226, 0.1208, and 0.0390 in the municipality order above.

The observed recovered and disrupted corridor travel times define disruption
penalties of 5, 70, 50, 30, and 40 minutes, respectively. Baseline failure
scores \(\Phi_\ell=(0.1350,0.4650,0.5438,0.3157,0.3759)\) are constructed
from observed delay ratios and restoration evidence; they are scenario
generators, not observed probabilities. The primary acute-access setting
uses \(\eta_r=0.25\), \(\delta_r=0\), a 60-minute timely-access threshold,
\(\underline\Phi_\ell=0.10\Phi_\ell\), \(\alpha_\ell=0.50\), and
\(\kappa=2\). The sector costs, budgets, capacity operational shares, and
throughput conversion are fixed before optimization and reported with the
replication files. Their role is comparative policy analysis, not
engineering cost estimation.

\subsection{Primary Noto results}
\label{sec:noto-results}

For each ambiguity endpoint \(\rho\in\{0,0.25\}\), the analysis scans all
\(5^5=3{,}125\) road-grid vectors in \eqref{eq:road-grid}; 996 satisfy the
road budget and are evaluated by the exact fixed-road oracle. Table
\ref{tab:noto-policy-results} compares the no-investment, optimized
no-road, and coordinated all-sector plans. Integrated mitigation reduces
modeled loss substantially relative to no investment. Road retrofitting also
has material incremental value after renovation and capacity are optimized:
96.858 modeled-loss units at \(\rho=0\) and 143.308 at \(\rho=0.25\).

\begin{table}[t]
\centering
\caption{Primary Noto policy comparison under capped TV ambiguity.}
\label{tab:noto-policy-results}
\small
\begin{tabular}{lrrrr}
\toprule
& \multicolumn{2}{c}{Worst-case modeled loss} &
\multicolumn{2}{c}{Reduction versus no investment}\\
\cmidrule(lr){2-3}\cmidrule(lr){4-5}
Policy & \(\rho=0\) & \(\rho=0.25\) & \(\rho=0\) & \(\rho=0.25\)\\
\midrule
No investment & 2{,}405.818 & 2{,}463.028 & 0.000 & 0.000\\
Renovation + capacity, \(y=0\) & 1{,}394.045 & 1{,}451.254
  & 1{,}011.774 & 1{,}011.774\\
All sectors & 1{,}297.187 & 1{,}307.946
  & 1{,}108.632 & 1{,}155.081\\
\midrule
Road value versus optimized no-road plan & 96.858 & 143.308 & -- & --\\
\bottomrule
\end{tabular}
\par\smallskip
\footnotesize\textit{Notes.} Lower is better. The outcome is modeled
time-sensitive people without timely access under the declared scenario;
it is not an estimate of earthquake deaths. Each all-sector row is
globally best over the declared five-level road grid.
\end{table}

At both endpoints, the selected road vector is
\begin{equation}
 \widehat y=(0,1,0.5,1,0),
\label{eq:noto-road-policy}
\end{equation}
in the corridor order stated above. The selected renovation plan is also
unchanged: it renovates 0.2769 of the Wajima zone and no other zone on the
declared grid/continuous renovation layer. Capacity has multiple
evaluation-equivalent allocations at \(\rho=0.25\). Consequently, a
different reported capacity vector at that endpoint is a zero-value tie,
not evidence of a substantively different capacity policy.

\subsection{What ambiguity changes, and what it does not}
\label{sec:noto-ambiguity}

Table~\ref{tab:noto-ambiguity} separates the ambiguity premium from policy
adaptation. At \(\rho=0.25\), ambiguity raises the value of the selected
plan by 10.759 modeled-loss units (0.823\% of the robust value), while the
complete-policy adjustment value in \eqref{eq:complete-adaptation} is zero.
The ambiguity set moves only 0.0499 probability mass despite the nominal
radius of 0.25 because the density cap restricts profitable transfers. The
observed premium is below the state-range audit bound in
\eqref{eq:ambiguity-premium-bound}.

\begin{table}[t]
\centering
\caption{Ambiguity and policy-identification diagnostics for the primary Noto grid.}
\label{tab:noto-ambiguity}
\small
\begin{tabular}{lrr}
\toprule
Diagnostic & \(\rho=0\) & \(\rho=0.25\)\\
\midrule
Selected-plan nominal value & 1{,}297.187 & 1{,}297.187\\
Selected-plan robust value & 1{,}297.187 & 1{,}307.946\\
Ambiguity premium & 0.000 & 10.759\\
Ambiguity premium as percent of robust value & 0.000\% & 0.823\%\\
Positive probability mass moved by nature & 0.0000 & 0.0499\\
TV state-range upper bound on premium & 0.000 & 11.415\\
\(\Delta_\rho^{\mathrm{all}}\) from \eqref{eq:complete-adaptation} & 0.000 & 0.000\\
Near-optimal portfolios within 0.01\% & 12 & 6\\
Near-optimal portfolios within 0.05\% & 47 & 19\\
Near-optimal portfolios within 0.10\% & 78 & 47\\
Near-optimal portfolios within 0.50\% & 93 & 93\\
\bottomrule
\end{tabular}
\par\smallskip
\footnotesize\textit{Notes.} Near-optimal counts are road portfolios;
for each portfolio, renovation and capacity are optimized by the fixed-road
oracle. Counts diagnose a flat discrete landscape and should not be
interpreted as distinct equally implementable policies without additional
operational screening.
\end{table}

The primary result is therefore deliberately narrow. Capped
decision-dependent ambiguity changes the worst-case evaluation and identifies
the amount of ambiguity premium, but it does not change the value-relevant
selected road or renovation plan in this Noto configuration. This is not a
failure of the solver: the conclusion follows from exhaustive grid
enumeration and complete-policy evaluation. Conversely, it would be
incorrect to claim that ambiguity changes the optimal policy solely because
an alternative tied capacity allocation is returned.

The common loss floor explains part of this stability. After renovation,
the selected plan has modeled loss exposure of 2,648.672 units, of which
1,286.428 units remain a state-common floor. That floor represents 99.17\%
of the robust objective at \(\rho=0\) and 98.35\% at \(\rho=0.25\).
Uncertainty therefore affects a comparatively small access-sensitive
component. The near-optimal counts in Table~\ref{tab:noto-ambiguity} show
that the policy landscape is also partly flat. These diagnostics are more
informative for practitioners than forcing a policy switch by selecting
ambiguity parameters ex post.

\subsection{Supplementary cross-case checks}
\label{sec:cross-case}

Turkey and Nepal are used to assess transferability of the empirical
workflow, not to claim that each dataset validates every model component.
The Turkey 2023 instance combines public Zenodo building-damage data,
HOTOSM/HDX destroyed-building and health-facility layers, WorldPop
population, and OpenStreetMap roads. Its road-failure scores, costs,
capacities, budgets, and ambiguity parameters are constructed or
scenario-calibrated. It is a static probability-only cross-check and hence
does not validate the Noto conditional-performance channel in
\eqref{eq:conditional-delay}.

The Nepal 2015 data use the public DrivenData Richter's Predictor building
damage labels. They provide real building-damage observations but anonymize
the geography, preventing a direct building-to-OSM road linkage. The Nepal
analysis is consequently an exposure/access stress test with constructed
corridors and a damage-rate proxy, not a geographically complete validation
of road failure.

\begin{table}[t]
\centering
\caption{Supplementary exact-grid cross-case checks at \(\rho=0.10\).}
\label{tab:cross-case-results}
\small
\resizebox{\textwidth}{!}{%
\begin{tabular}{p{2.2cm}rrrrp{4.7cm}}
\toprule
Case & Feasible grid portfolios & Best robust value & No-investment value &
Road value versus optimized no-road plan & Exact-grid finding\\
\midrule
Turkey 2023 &
347 & 106{,}451.774 & 222{,}162.759 & 138.163 &
The same \(y=(0,1,0,1,0.75)\) is selected over
\(\rho\in\{0,0.05,\ldots,0.25\}\); \(\Delta_\rho^{\mathrm{all}}=0\).\\
Nepal 2015 access stress test &
983 & 195{,}234.523 & 237{,}371.955 & 6{,}897.940 &
The same \(y=(0,1,1,0,0)\) is selected over
\(\rho\in\{0,0.05,\ldots,0.25\}\); \(\Delta_\rho^{\mathrm{all}}=0\).\\
\bottomrule
\end{tabular}}
\par\smallskip
\footnotesize\textit{Notes.} Values are modeled-loss units and are
comparable only within a case. In both cases every feasible five-level road
portfolio is evaluated exactly by Algorithm~\ref{alg:fixed-y}; the result is
not a continuous-\(y\) certificate.
\end{table}

The three cases show that the workflow can produce nonzero road-sector value
without requiring that ambiguity change the selected long-lived road plan.
This is a practically useful outcome: under the declared assumptions, the
selected plan is a no-regret choice across the tested ambiguity radii. It
does not imply a universal rule such as \emph{always harden the cheapest
link} or \emph{always design for the worst state}. The selected portfolios
depend on observed exposure, access bottlenecks, route penalties, residual
risk, capacity, budgets, and the stated uncertainty envelope.

\subsection{Sensitivity, solution scope, and limitations}
\label{sec:results-limitations}

The capped-TV specification is varied over \(\rho\), and the density cap,
residual floor, road budget, cost coefficients, response threshold, and
moment-band tolerances are suitable sensitivity dimensions. They should be
declared before interpreting policy changes. The optional moment-aware
extension in \eqref{eq:moment-marginal}--\eqref{eq:moment-second} is useful
when the analyst has a defensible hazard ensemble or wishes to bound changes
in link marginals and dependence explicitly. In the present scarcity
sensitivity it can alter a short-run capacity allocation while the
long-lived road and renovation selections remain stable; its tolerances are
therefore reported as scenario sensitivity, not as calibrated uncertainty
truth.

The following scope statements apply to all reported results:
\begin{enumerate}
 \item The observed data calibrate population, a damage proxy, facility
       locations or beds, and road/access layers where available. They do
       not validate the entire optimization model or establish causal
       retrofit effects.
 \item \(\Phi_\ell\), \(\underline\Phi_\ell\), \(\alpha_\ell\), costs,
       budgets, capacity operational shares, \(\rho\), \(\kappa\), and
       moment tolerances are constructed or scenario-calibrated. They must
       not be described as directly observed.
 \item The five-link policies are exhaustive-grid optima. They are not
       continuous global optima, and the weak five-link SBB lower bound in
       Table~\ref{tab:sbb-root-diagnostic} cannot support such a claim.
 \item The outcome is a declared modeled acute-access loss. It is not a
       validated forecast of deaths, a clinical survival model, or an
       estimate of the causal impact of a specific retrofit project.
 \item Policy stability is an empirical finding conditional on the declared
       grid and parameters. A policy switch should be claimed only after a
       complete-plan evaluation gives
       \(\Delta_\rho^{\mathrm{all}}>0\) and the switch survives an adequate
       grid-resolution or continuous verification check.
\end{enumerate}

Accordingly, the empirical contribution is a reproducible, auditable
decision-analysis workflow: it combines observed damage and access evidence
with transparent scenario assumptions, quantifies the incremental value of
coordinated mitigation, and distinguishes a robust evaluation effect from a
decision-changing ambiguity effect. This distinction allows practitioners to
use stable policies as no-regret plans while identifying exactly which
assumptions would need better data or a finer decision grid before claiming
that uncertainty changes the recommended investment.
